%Abstract will go here.  Make sure it remains within the 350 word limit.
%\par To get a new paragraph
\par This dissertation studies the problem of identifying disease clusters
associated with the highest mortality rates among patients in
electronic health record databases. We introduce the maximum
mortality rate clique problem, a fractional combinatorial
optimization problem defined on a comorbidity graph whose vertices are
disease categories and whose edges reflect statistically significant
pairwise co-occurrence. The problem seeks a frequent clique, supported
by a sufficiently large shared patient population, that maximizes the
mortality rate of the patients simultaneously diagnosed with all
diseases in the clique.
\par We establish that the mortality rate function is neither additive nor
monotonic, ruling out simple greedy strategies, and prove
NP-completeness of the core decision problem and several related
variants. These results motivate the development of exact solution
methods of increasing sophistication.
\par We develop two exact combinatorial algorithms: a modified
Bron--Kerbosch enumerative algorithm and a combinatorial
branch-and-bound algorithm that incorporates an upper-bounding scheme
to prune the search tree aggressively. We present two mixed-integer
linear programming formulations based on auxiliary product-variable
linearization and the Charnes--Cooper transformation. These
formulations become intractable as the patient population grows. To
overcome this, we introduce a logic-based Benders decomposition
reformulation that eliminates all patient-indexed variables and embeds
naturally in a decomposition branch-and-cut framework. A hybrid
variant further integrates the combinatorial branch-and-bound
algorithm to generate stronger valid inequalities, substantially
improving solution times and optimality gaps over the pure
logic-based Benders decomposition approach.
\par We further extend the framework to incorporate side constraints, a
size-restricted variant, and a marginal mortality rate variant. We
evaluate all proposed methods using the real-world electronic health
record dataset MIMIC-IV.